\section{Actual physical-source charts and the full Schwartz obstruction}
\label{sec:full-schwartz-cutoff-obstruction}

\paragraph{Status.}
The complete post-writing author reading and two independent
whole-proof readings have passed.  Their source and final-delta
checks retain their exact separately recorded scope.

The compact-frequency charts in MA-PI-005--008 can be connected
to the actual finite source vectors.  A separate finite-jet
argument then removes the restriction to those charts within
the stated archimedean Schwartz family.  This step does not
identify the global Schwartz source with the coordinate-flat
local source.

Use one of the following three actual source families.
In the rational principal case put
\[
 \mathscr X_0=\{f\in\mathcal S(\mathbb R):f(-t)=f(t),
                         \ f(0)=0,\ \int f=0\}.
 \tag{MA-PI.190}
\]
Its finite source is spherical, and $S=\{\infty\}$.
For a nonprincipal primitive Dirichlet character $\chi$ of
conductor $q$, let $\mathscr X_\chi$ be the Schwartz parity
space of parity $\chi(-1)$ and take
\[
 S=\{\infty\}\cup\{p:p\mid q\},\qquad
 \varphi_f=f\otimes\prod_{p\mid q}
                  \kappa_p\mathbf1_{\mathbb Z_p^\times}.
 \tag{MA-PI.191}
\]
The outside-$S$ vector is $\mathbf1_{\mathbb Z_p}$, and
$\kappa_p$ is the actual inverse unit character in MA-CQ-001.
Finally, for $K=\mathbb Q(\zeta_3)$ and the non-induced
$\psi=\psi_{269,3}$, let $\mathscr X_\psi$ be radial
Schwartz space on $\mathbb C$.  Use exactly
$S=\{\infty,\mathfrak p,\mathfrak q\}$ and
\[
 \varphi_f=f\otimes h_{\mathfrak p}\otimes\kappa_{\mathfrak q},
 \qquad q=269^2,
 \tag{MA-PI.192}
\]
with the finite factors of (MA-PI.120)--(MA-PI.121).
In both nonprincipal families the ramified finite factor
removes both global polar terms for every archimedean $f$.
All three domains have their actual induced Schwartz
Fr\'echet topology.

Write $d=1$ in the rational cases and $d=2$ in the complex
case.  The normalized physical dilation and centered
translation are
\[
 (D_yf)(z)=e^{-y/2}f(e^{-y/d}z),\qquad
 (\tau_yh)(x)=h(x-y).
 \tag{MA-PI.193}
\]
The actual centered theta map satisfies
$\Theta(D_yf)=\tau_y\Theta(f)$ in each family.  This follows
directly from (MA-PI.70), (MA-CGP.2), and the fixed finite
translation filter (MA-PI.129), respectively.

Fix either one of the actual Hilbert metrics
\[
 \begin{aligned}
 H^{\rm sym}_{a,s}&=L^2(e^{2a|x|}(1+x^2)^{s/2}dx),
       &&0<a<\tfrac12,\ s\ge0,\\
 H^{\rm one}_{b,s}&=L^2(e^{2bx}(1+x^2)^{s/2}dx),
       &&|b|<\tfrac12,\ s\ge0.
 \end{aligned}
 \tag{MA-PI.194}
\]
Let $P_Y$ be the orthogonal projector onto the closure of
the full original global PW source image at physical level
$Y$, in that fixed metric.  It is not the projection only
onto the particular finite-vector family above.
The physical cutoff is
$\mathfrak B_{S,Y}=\mathcal F_S^{-1}
M_{\rho_Y(|\xi|_S)}\mathcal F_S$, with
$\rho_Y=1$ on $[0,Y/2]$ and $\rho_Y=0$ on $[2Y,\infty)$.

\paragraph{MA-PI-009 (no nonzero full-Schwartz transport on these actual families).}
Let $\mathscr X$ be any one of (MA-PI.190)--(MA-PI.192),
and let $\mathscr I_\eta^c$ be its actual centered smooth
theta ideal.  Suppose
$T:\mathscr X\to\mathscr I_\eta^c$ is linear and continuous into
the full exponential-Schwartz topology and satisfies
$TD_y=\tau_yT$.
Assume there is one $Y_0$ such that, for every $Y\ge Y_0$
and every source vector in this family, the exact plateau
identities hold:
\[
 \begin{array}{ll}
 P_YTf=Tf,&\mathfrak B_{S,Y}\varphi_f=\varphi_f,\\
 P_YTf=0,&\mathfrak B_{S,Y}\varphi_f=0.
 \end{array}
 \tag{MA-PI.195}
\]
Then $T=0$ on the entire stated Schwartz family.
Any exact cutoff intertwining on a larger source domain
necessarily has these plateau consequences on its restriction.
Only these consequences are assumed; an arbitrary transition
output need not stay in the same fixed finite-vector family.

For the principal rational family there is also a weaker-target
conclusion in the exact printed central norm:
$T:\mathscr X_0\to J_s^{\rm pr}$ may be merely Hilbert-valued
and continuous when $s>1$.  With the same equivariance and
the printed $A=1/2$ orthogonal projections, (MA-PI.195)
again implies $T=0$.
This additional case explicitly uses the critical-line subset
theorem in MA-PI-008.  The smooth cases instead use the full
classical zero count with multiplicities from MA-PI-007,
without a narrower zero-location assertion, A, or M.

\begin{proof}
\emph{1. The actual compact-frequency charts.}
For $v\in\mathcal D(\mathbb R)$ define on the positive
real frequency ray
\[
 u_v(t)=t^{-1/2}v(-\log t),\quad t>0,
 \qquad f_v=\mathcal F_\infty^{-1}u_v.
 \tag{MA-PI.196}
\]
In the real cases extend by the required parity.
In the complex case interpret the first formula as a norm
profile: $u_v(z)=|z|^{-1}v(-\log|z|^2)$ for $z\ne0$.
These are smooth functions with compact support away from
the additive origin; their inverse Fourier transforms are
genuine Schwartz inputs.  On each fixed compact support
in the logarithmic coordinate the map to Schwartz space is
continuous, hence it is continuous on the strict LF domain.

For the principal rational sector, $u_v(0)=0$ already,
and
\[
 \int_{\mathbb R}u_v(\xi)d\xi
       =2\int_{\mathbb R}e^{-x/2}v(x)\,dx.
 \tag{MA-PI.197}
\]
Thus $f_v$ belongs to $\mathscr X_0$ exactly for
$v\in\mathcal D_0$.  No moment restriction is required
in either nonprincipal sector, since its finite factor
already removes both polar terms.

Replacing $v$ by $\tau_yv$ gives
$u_{\tau_yv}(z)=e^{y/2}u_v(e^{y/d}z)$.
The real Jacobian of this frequency change is $e^{-y}$,
in both dimensions, so
\[
 f_{\tau_yv}(z)=e^{-y/2}f_v(e^{-y/d}z)=D_yf_v(z).
 \tag{MA-PI.198}
\]
Consequently $v\mapsto Tf_v$ has exactly the translation
equivariance required in MA-PI-007, or in MA-PI-008 for
the principal Hilbert-valued case.

\emph{2. All finite Fourier shells and the plateau separation.}
In the principal sector the finite norm factor is one.
For a primitive Dirichlet character, (MA-CQ.5) places
each ramified finite transform on the single shell
$|\xi_p|_p=p^{f_p}$, because the rational different
exponent is zero.  Their product is exactly $q$.
The finite phases do not change the norm support.
In the cubic case (MA-PI.125)--(MA-PI.127) give the two
distinct product shells $9q$ and $27q$; both are retained.
Accordingly set
\[
 (c_{\min},c_{\max})=
 \begin{cases}
 (1,1),&\eta=\mathbf1\text{ over }\mathbb Q,\\
 (q,q),&\eta=\chi\text{ primitive nonprincipal},\\
 (9q,27q),&\eta=\psi_{269,3}\text{ over }K.
 \end{cases}
 \tag{MA-PI.199}
\]
The infinite-place norm on the Fourier support is
$e^{-x}$.  Write $c=-\log Y$.  If
\[
 \begin{array}{ll}
 \operatorname{supp}v\subset[c+\log(2c_{\max}),\infty),
       &\mathfrak B_{S,Y}\varphi_{f_v}=\varphi_{f_v},\\
 \operatorname{supp}v\subset(-\infty,c+\log(c_{\min}/2)],
       &\mathfrak B_{S,Y}\varphi_{f_v}=0,
 \end{array}
 \tag{MA-PI.200}
\]
because the physical cutoff is then constant on every
finite Fourier shell.  This is a literal product-norm
calculation, without replacing the two cubic shells by
one new conductor or changing the physical level.

Two ordered compact logarithmic supports separated by
more than
\[
 R=\log(4c_{\max}/c_{\min})
     =\begin{cases}\log4,&\mathbb Q,\\\log12,&K,\psi_{269,3},\end{cases}
 \tag{MA-PI.201}
\]
admit a $c$ for which the left support is in the zero
plateau and the right support is in the one plateau.
Translate both inputs by $\tau_{-L}$ and replace $c$
by $c-L$.  The corresponding level is $Y_L=e^LY$,
so it exceeds the fixed $Y_0$ for all sufficiently
large $L$.  Equation (MA-PI.195) makes one output
orthogonal to the other, since the same orthogonal
$P_{Y_L}$ fixes one and kills the other.
MA-PI-007 now gives $Tf_v=0$ for all $v$ in the
specified chart.  In the printed principal Hilbert case
MA-PI-008 gives the identical conclusion on $\mathcal D_0$.

\emph{3. A finite-jet lemma in the actual target norm.}
We give the needed vector-valued form explicitly.
Let $\mathscr H$ be either fixed Hilbert space in
(MA-PI.194), or $H_s^{\rm pr}$, and suppose a continuous
linear map $R:\mathcal S(\mathbb R^d)\to\mathscr H$
vanishes on every compact smooth function supported
away from the origin.  Continuity supplies one integer
$N$ and one finite Schwartz-seminorm bound for $\|Ru\|$.
If all derivatives of $u$ through order $N$ vanish at
zero, Taylor's formula and a fixed $\chi\in C_c^\infty$
equal
to one near zero give
\[
 \sum_{|\alpha|\le N}\sup_z
 (1+|z|)^N\bigl|\partial^\alpha
                    (\chi(z/\varepsilon)u(z))\bigr|=o(1).
 \tag{MA-PI.202}
\]
The complement $(1-\chi(z/\varepsilon))u$ vanishes
near zero and is a Schwartz limit of compact functions
with the same property.  Its image under $R$ is zero.
The continuity bound and (MA-PI.202) therefore give $Ru=0$.
Choosing a finite jet basis yields vectors $h_\alpha$
in the actual Hilbert space such that
\[
                   Ru=\sum_{|\alpha|\le N}
                                 h_\alpha\partial^\alpha u(0).
 \tag{MA-PI.203}
\]
The argument uses one target norm and one finite order;
there is no unsupported density assertion in the
whole nonflat Schwartz space.

If in addition
$R(e^{y/2}u(e^{y/d}\cdot))=\tau_yRu$, comparison of
independently prescribed jets gives
\[
 \tau_yh_\alpha=e^{(1/2+|\alpha|/d)y}h_\alpha.
 \tag{MA-PI.204}
\]
A translation eigenvector with this real exponent is,
as a distribution, a constant multiple of
$e^{-(1/2+|\alpha|/d)x}$: differentiate the identity
with respect to $y$ in distributions and solve the
resulting first-order equation.
No nonzero such exponential belongs to any of the
whole-line Hilbert spaces under consideration.
In a one-sided weight, it grows at one end unless the
exponents cancel, in which case its squared weighted
norm is the divergent integral of $(1+x^2)^{s/2}$.
The symmetric weight also forbids integrability at both
ends.  Hence every $h_\alpha$ is zero.

The lemma applies to the real parity and complex radial
spaces by first composing with their continuous compact
averaging projector on Schwartz space.  This projector
preserves support away from zero and commutes with the
normalized dilations.  Thus it supplies an equivariant
extension to the full real-coordinate Schwartz space
without altering the given map on its sector.
In the nonprincipal applications the map in this lemma is
explicitly $R(u)=T\mathcal F_\infty^{-1}(P_{\rm sector}u)$.
Step~2 supplies its compact-frequency vanishing, and
(MA-PI.198) supplies its normalized dilation covariance.

\emph{4. The principal two-moment variant.}
In Fourier coordinates the principal domain is
$\mathscr Z_0=\{u\in\mathcal S(\mathbb R)^{\rm even}:
u(0)=0,\int u=0\}$.
It is a closed finite-codimensional subspace with a
continuous projection from even Schwartz space,
obtained by correcting the value and integral with
two fixed even test functions.  Extend the map
$u\mapsto T\mathcal F^{-1}u$ by that projection,
and call the extension $R$.
On compact even tests away from zero, Step~2 says
that $R$ vanishes on the kernel of the integral.
Choose one such test $v_*$ with integral one, and put
$h_*=Rv_*$.  Then $R-h_*\int$ vanishes on all those
away-from-zero tests.  The finite-jet lemma, after even
projection, shows that on $\mathscr Z_0$ the original
map has a finite expansion
\[
       T\mathcal F^{-1}u=\sum_{\substack{2\le j\le N\\j\ \text{even}}}
                                     h_j u^{(j)}(0).
 \tag{MA-PI.205}
\]
Here the zeroth jet and integral vanish by definition;
odd jets vanish by parity.  The remaining finite jets
can be prescribed independently in $\mathscr Z_0$:
choose a local even jet polynomial times a cutoff,
then correct its integral by a remote even bump.
The original, unextended map is equivariant on this
invariant domain.  Comparing these independent jets
therefore gives $\tau_yh_j=e^{(j+1/2)y}h_j$.
The same absence of Hilbert translation eigenvectors
forces every coefficient to vanish.

For the smooth-target cases, composing the given map
with the continuous injective inclusion
$\mathcal E\to\mathscr H$ provides exactly the target
norm used above.  Its vanishing in that norm is its
vanishing as a smooth map.  For the printed Hilbert
case no additional target inclusion is needed.
Steps~3--4 prove the asserted zero conclusion on each
entire Schwartz family.
\end{proof}

\paragraph{Exact source and remaining domain boundaries.}
The physical source charts above are genuine global PW
inputs.  The rational principal source is Masullo,
Definition~7.3.2.3 and (7.3.2.12)--(7.3.2.17),
pp.~354--355; its two polar conditions are retained.
The nonprincipal finite vectors, Gauss shells and
compact-character projections are the already checked
MA-CQ-001 and MA-CGP-001 constructions.  The cubic
different place and both actual Fourier shells are
those of MA-PI-004.  The physical frequency ball is
p.~516, Step~II, with the thin-transition plateaus
used on p.~519.  Orthogonal PW projections are
Definition~11.2.1.3, pp.~488--489.

For the printed norm, the unitary is exactly
$(U_bh)(\nu(e^x)k)=m^{-1/2}\eta(k)^{-1}e^{bx}h(x)$,
with the actual compact Haar mass $m$ and $b=A-1/2$,
as proved in MA-PI-002 and (MA-PI.148).
Thus the centered equivariance corresponds to
$\mathfrak T D_y=e^{by}\tau_y\mathfrak T$ in printed
norm coordinates.  No orthogonal projector is moved
by a mere norm equivalence.

The theorem treats the stated Schwartz families with
their stated continuity, ideal range and all-level
plateau hypotheses.  It does not say that compact
frequency functions are dense in every Schwartz
source; Step~3 supplies the missing finite-jet argument.
It does not extend a map known only on the coordinate-flat
local LF source, construct an inverse for an arbitrary
global-to-local map, or identify every possible
level-dependent or Gram-corrected comparison with
the class just excluded.  Those are distinct interfaces.

In the inherited cote~1 BR-TR-01/02 placement,
scans079--080 and111--112, the obstruction is to an
actual continuous comparison before trace cyclicity.
The principal, Dirichlet and cubic sectors retain
their same finite Artin local systems in SGA~1 V~\S7
and Deligne, Definition~6.1 and \S\S6.2--6.3.
The cubic $\mathfrak p$ filter changes the test vector,
not the primitive conductor or motive.
No new purity or faithful full-jet positivity is inferred,
and the independent FT002 arithmetic gate remains open.
